%% Resumo
\begin{resumo}
O presente trabalho descreve o desenvolvimento do \textit{SyncCarreira}, uma plataforma \textit{Software as a Service} (\textit{SaaS}) voltada para a automação e otimização dos processos de orientação vocacional em instituições de ensino. Motivada pela ineficiência dos métodos manuais, que sobrecarregam os psicólogos escolares e dificultam o acompanhamento longitudinal dos estudantes, a solução centraliza dados, automatiza a correção de testes e fornece \textit{dashboards} analíticos para a gestão de turmas. Utilizando tecnologias como \textit{React} no \textit{front-end}, Java com \textit{Spring Boot} no \textit{back-end} e MySQL para persistência, a aplicação foi estruturada com base em uma arquitetura em camadas e metodologias ágeis Scrum. Os resultados esperados incluem a redução da burocracia, maior engajamento dos alunos e suporte técnico qualificado para a tomada de decisão vocacional, garantindo conformidade com a Lei Geral de Proteção de Dados (LGPD).

\vspace{\onelineskip}
\noindent
\textbf{Palavras-chave}: Orientação Vocacional. Automação. SaaS. Tecnologia Educacional. Gestão de Dados.
\end{resumo}
